\section{Constructive Recovery and Verification}
\label{sec:constructive_recovery}

This section suppresses the fixed update epoch. It supplies the locality
criterion, constructive reference search, and worked instance used in the
main paper. All certificates remain conditional on the received wind table.

\subsection{Locality, Effective Freedom, and Collapse}
\label{subsec:locality_resolution}

Let $\mathfrak C$ be the connected components of
$\mathcal G^{\mathrm{task}}$ and write $S=\mathcal V^{\mathrm{aff}}$.
Reachability gives the exact identity
\begin{equation}
\mathcal V^R=M^A\cap\bigcup_{C\in\mathfrak C:\,C\cap S\ne\varnothing}C.
\label{eq:supp_component_locality}
\end{equation}
In particular, $\mathcal V^R\subsetneq M^A$ if and only if there is a
mutable assignment whose component contains no seed. A sufficient way to
establish this is a partition of the task set such that every declared
support $T(g)$ lies within one block, seeds lie in only some blocks, and
an unseeded block contains a mutable task. This is a statement about the
declared supports; it is not inferred from geographic distance or incumbent
constraint slack.

Define
\begin{equation}
E^A=\{i\in M^A:\exists\Pi\in\mathcal Y,\ a_i\ne\bar a_i\}.
\label{eq:supp_effective_assignments}
\end{equation}
Nominally mutable coordinates outside $E^A$ are already fixed by structural
consistency or other commitments. Determining $E^A$ can itself require
search; the procedure need not compute it to obtain a safe neighborhood.

\medskip
\noindent\textbf{Proposition 4 (strictness and collapse).}
Under the fixed support and common-domain conventions,
\begin{equation}
\mathcal X_3\subsetneq\mathcal X_4
\quad\Longleftrightarrow\quad E^A\setminus\mathcal V^R\ne\varnothing.
\label{eq:supp_strict_spaces}
\end{equation}
If a reached support contains $M^A$, then $\mathcal X_3=\mathcal X_4$.
Moreover, $\mathcal X_2\subsetneq\mathcal X_3$ holds exactly when some
plan in $\mathcal X_3$ changes an assignment from the incumbent.
\emph{Proof.}
The only additional restrictions in $\mathcal X_3$ relative to
$\mathcal X_4=\mathcal Y$ are the assignment equalities outside
$\mathcal V^R$. A plan violating any one exists precisely under the
condition in \eqref{eq:supp_strict_spaces}. A reached support containing
$M^A$ puts all mutable assignments into the closure, leaving no such
equalities. Finally, $\mathcal X_2$ fixes all assignments and otherwise
has the same mutable sequence, path, and transfer releases as
$\mathcal X_3$. \hfill$\square$

Strictness of $\mathcal X_3\subsetneq\mathcal X_4$ does not imply
$\mathcal F_3\subsetneq\mathcal F_4$: the extra structural plans could
all be infeasible. Likewise, a proper neighborhood can leave identical
spaces when its excluded assignments are structurally fixed.

For example, if $\mathcal T_u=\mathcal V^{\mathrm{rem}}$, any constraint
depending on a mutable route coordinate of $u$ has full task incidence
under the stated attribution. A nonempty seed set then reaches every task.
Shared stock and finite bay capacity can also join otherwise separated
groups. Such edges are retained. Constraint elimination is permissible
only with a proof of satisfaction throughout the common domain; removing
an incumbent-slack constraint is insufficient. Smaller exact supports may
be proved from fixed endpoints or other structural restrictions.

For searching, merge levels only after equality of their decision spaces
in the searched representation is certified. Equal release masks suffice;
the implication $E^A\subseteq\mathcal V^R$ also suffices for Levels 3--4
in the full domain. Without an equality proof, retain both searches.
An equivalence class keeps its smallest original label. Hence, when
$\mathcal X_3=\mathcal X_4$, successful reassignment has label 3 and never
a separate minimum of 4. This preserves the original definition of the
minimum rather than renumbering levels.

\subsection{Declared Reference Representation}
\label{subsec:reference_domain}

The reference procedure specializes the fixed operation protocol to
durations and ground energies that depend on the UAV, site, slot, and
observed residual mode, but not on decision quantities, entry time, or
entering battery. Depot preparation restores full charge in its prescribed
turnaround time. This is an explicit simulation convention, not a model
of charging electrochemistry. It is chosen before the paired evaluations
and held fixed across all levels. Quantity-dependent service/charging
durations require a different solver and are outside the linear reference
reduction.

Let $\mathfrak R$ specify finite libraries
$\mathscr P_u^{\mathfrak R}(i,j)$ and
$\mathscr W_u^{\mathfrak R}(i)$ for nominal endpoint pairs and returns
from a site or observed origin to the depot. A concrete construction uses
a finite directed waypoint graph whose edges are straight segments in
$\mathcal W$, and enumerates all walks of at most $H_P$ edges.
Repeated vertices are allowed: excluding all loops would be an additional
restriction in time-varying wind. Include the constant depot path and
append every incumbent/committed nominal curve and each supplied
pre-update reserve witness, even if it exceeds $H_P$.
Canonicalize geometrically identical arclength curves.
Empty endpoint libraries are explicit and cannot be filled differently
at different levels.

The waypoint graph includes all task locations, the depot, and each
observed origin; any extra vertices, $H_P$, and library ordering are fixed
before comparing levels using only received information. Libraries contain
geometry, not precomputed departure-time costs. Every evaluation splits
paths at spatial boundaries and integrates the appropriate time bins at
the actual induced departure time. The support rules and the resulting
$\mathcal V^R$ are also fixed across levels. Supports valid on the original
$\mathcal Y$ remain conservative on the restricted domain.

Let $\mathcal Y^{\mathfrak R}$ impose the nominal libraries in addition
to the original structural and commitment rules, with continuous transfer
variables. Write $\mathcal F_k^{\mathfrak R}$ for plans in
$\mathcal X_k\cap\mathcal Y^{\mathfrak R}$ satisfying all nominal
conditions and
\begin{equation}
\mathscr W_u(z)\cap\mathscr W_u^{\mathfrak R}(\mathbf x_z)
\ne\varnothing\quad\forall (u,z)\in\mathcal Z^{\mathrm{res}}(\Pi).
\label{eq:represented_reserve}
\end{equation}
The shorthand library argument $\mathbf x_z$ denotes the matching site or
origin. A reserve witness is not an intervention coordinate: it may be
chosen at every level, including Level-0.
Then
\begin{equation}
\mathcal F_k^{\mathfrak R}\subseteq\mathcal F_k,
\qquad k^F\le k_{\mathfrak R}^F
:=\min\{k:\mathcal F_k^{\mathfrak R}\ne\varnothing\}.
\label{eq:reference_scope_relation}
\end{equation}
The incumbent and its pre-update reserve witnesses are validated in
$\mathfrak R$ before the instance is admitted to the reference benchmark.

Enlarging nominal or reserve libraries while retaining the same support
graph, operation protocol, and hierarchy can only enlarge each feasible
set and decrease $k_{\mathfrak R}^F$. This gives a useful representation
check. Rebuilding a different hierarchy during refinement invalidates that
particular comparison. Finite-library exhaustion gives no emptiness proof
for arbitrary piecewise-linear paths.

\subsection{Reduction to Linear Feasibility}
\label{subsec:linear_reference_search}

Enumerate assignment vectors, active-prefix sequences, and nominal
library paths satisfying structural compatibility, commitments, and the
level's unreleased coordinates. Call a resulting skeleton
$\sigma=(\mathcal A,\mathcal S,\mathcal P)$. Assignments may equivalently
be derived from a consistent sequence, avoiding duplicate enumeration.
There are finitely many skeletons because the slot bounds and libraries
are finite.

For a skeleton, perform exact wind integration and the prescribed ground
operations. These times do not depend on loading or delivery quantities.
A late nominal operation, fixed ground completion failure, or bay conflict
excludes the whole skeleton. Bay intervals are checked with departures
processed before arrivals at equal times, consistently with the half-open
interval convention; zero-duration intervals occupy no bay.

Let $\mathbf l$ contain all remaining delivery and loading variables.
The payload recursion gives affine functions
$q_{u,h}^{\pm}(\mathbf l)$. Define the known coefficient
$c_{u,h}=\gamma_u v_u\tau_{u,h}$. Then
\begin{equation}
E_{u,h}^{\mathrm{fly}}(\mathbf l)
=c_{u,h}\bigl(m_u^0+m q_{u,h-1}^{+}(\mathbf l)\bigr).
\label{eq:supp_affine_energy}
\end{equation}
Battery states are affine by induction, including the constant reset to
$B_u$ at depot departure. Fixed service energies preserve this property.
For each reserve-check state, enumerate one path from its finite return
library. Its start time and duration are already fixed by $\sigma$;
its payload and battery remain affine in $\mathbf l$.
Denote the vector of reserve-path choices by $\omega$.
A return path finishing after $T$ excludes that choice, not all choices.

For each pair $(\sigma,\omega)$ retain the following constraints:
nonnegative transfers and payloads; delivery availability; payload
capacity; stock; per-task residual mandatory and demand bounds; battery
bounds; pre-recharge arrival and terminal reserves; and every selected
return-energy inequality. Add site-specific transfer zeros, padding,
commitment equalities, and the level's fixed transfers. If a relief visit
has fixed completion time above its deadline, impose $x_{u,h}=0$.
Thus the positive-delivery implication introduces no strict inequality
or additional binary variable within a skeleton.

The branch is a continuous linear feasibility problem
\begin{equation}
A\mathbf l\le b,\qquad C\mathbf l=d.
\label{eq:supp_branch_lp}
\end{equation}
Each allowed transfer can additionally be bounded by
$0\le x_{u,h},y_{u,h}\le C_u^{\mathrm{pay}}/m$.
These upper bounds are implied by delivery availability or depot payload
capacity (loading and unloading do not occur together). Fixed impossible
quantities reject the branch. Keeping the explicit bounds makes the
transfer polytope compact.

\medskip
\noindent\textbf{Proposition 5 (finite reference completeness).}
Under the prescribed transfer-independent ground protocol, a represented
plan is feasible at level $k$ if and only if at least one enumerated
branch has a feasible transfer vector.
\emph{Proof.}
A feasible plan determines an enumerated nominal skeleton. At each of its
finitely many reserve states, represented feasibility supplies at least
one return path, so an enumerated $\omega$ contains those witnesses.
Its transfer vector satisfies every row of \eqref{eq:supp_branch_lp}.
Conversely, a feasible vector for a retained branch has structurally
admissible nominal decisions, complete time records, every resource and
service condition, and explicit reserve witnesses. These are precisely
the represented feasibility requirements. \hfill$\square$

This construction does not assume independent UAV resource subproblems:
the LP retains all shared-stock rows, and the skeleton schedule retains
the full bay check. Its cost can be exponential in route, path, and reserve
choices. It is a small-instance correctness reference, with no real-time
or fleet-scale efficiency claim.

\subsection{Certificates, Arithmetic, and Search Budgets}
\label{subsec:recovery_certification}

A positive branch certificate contains $\sigma$, $\omega$, $\mathbf l$,
and recomputed induced states. An independent validator checks the
original predicates and level membership, not only the LP residuals.
For \eqref{eq:supp_branch_lp}, a sufficient exclusion certificate is
\begin{equation}
\begin{split}
\lambda&\ge0,\qquad A^{\mathsf T}\lambda+C^{\mathsf T}\nu=0,\\
b^{\mathsf T}\lambda+d^{\mathsf T}\nu&<0,
\end{split}
\label{eq:branch_exclusion_certificate}
\end{equation}
with unrestricted equality multipliers $\nu$. If a feasible vector
existed, premultiplication would give
$0\le b^{\mathsf T}\lambda+d^{\mathsf T}\nu$, a contradiction.
Equality rows can instead be encoded as pairs of inequalities.
Geometric/time rejection and certified structural inconsistency are also
valid exclusions and must carry their reason.

For a fully specified exact reference option, use rational primitive
geometry, wind, vehicle, and fixed-operation data. Cell intersections,
segment lengths, wind speeds, and bin-integrated times are then real
algebraic numbers. Exact algebraic arithmetic and sign decisions make the
branch tests well-defined. Encode all equalities and bounds as inequalities.
With $p$ free transfer variables, enumerate sets of $p$ linearly independent
active rows, solve their systems, and check the resulting points against
all rows. A nonempty bounded polyhedron has a vertex; consequently,
exhausting all such bases without a feasible point certifies emptiness.
When $p=0$, check the constant rows directly. This is finite but expensive.
Rational collinear-wind examples below require only rational arithmetic.

Floating-point LP software can propose witnesses and dual certificates,
but proof labels require validated arithmetic for the declared coefficients.
An exact check, or certified enclosures that establish every required
inequality and equality, may validate them. Unresolved signs, equality
residuals, or solver failure yield unresolved search. Rounded coefficients
and feasibility tolerances must not silently redefine the represented
mission. Equality and horizon-boundary cases require exact treatment or an
explicitly declared alternative numerical model.

For each level, fix a branch ordering and a branch/LP budget $B_k$;
also fix a total budget. A finite wall-clock limit may be imposed, with
its runtime recorded separately. Exhaustive search returns a witness
immediately after validation, an empty status only after all branches have
sound exclusion records, and unresolved otherwise. Proven infeasible
branches can be pruned; a missed or numerically unresolved branch prevents
an emptiness certificate. Complete exact search terminates for the finite
representation; bounded search need not decide it.

The outer procedure in the main paper searches increasing original labels.
Empty certificates are credited to every proven-equivalent label. Nesting
also excludes all smaller labels. A timeout at a smaller level does not
prevent searching a larger level for a usable witness if budget remains.
It does prevent claiming that witness's label is minimal without further
exclusion. The output contains
\begin{equation}
(\mathfrak R,\Pi^{\mathrm{wit}},\omega^{\mathrm{wit}},W,D,\ell,u,
\text{unresolved labels},\text{search counts}),
\label{eq:recovery_record}
\end{equation}
with absent witnesses represented explicitly. Bounds are
\begin{equation}
u=\min W,\qquad
\ell=\begin{cases}
+\infty,&4\in D,\\
1+\max(\{-1\}\cup D),&\text{otherwise}.
\end{cases}
\label{eq:supp_search_bounds}
\end{equation}
The bounds apply to $k_{\mathfrak R}^F$. A finite validated upper bound also
applies to $k^F$; the represented lower bound generally does not.
Use four explicit outcomes: minimum certified, witness with unresolved
minimum, no represented continuation, and unresolved without a witness.
Only an independent exclusion over the unrestricted domain can upgrade
the third outcome to unrestricted impossibility.

\balance
\subsection{Worked Locality and Recovery Instance}
\label{subsec:worked_recovery_instance}

All values use consistent normalized units. Set $t=0$, $T=4$,
$\mathcal W=[-1,1]^2$, and $\mathbf p_0=(0,0)$. Distinct recipients
1 and 2 are at $(1,0)$; recipients 3 and 4 are at $(-1,0)$.
Each requires exactly one commodity unit: $d_i=d_i^{\mathrm{mand}}=1$,
and no delivery has occurred. Set $m=1$, $S^{\mathrm{dep}}=4$,
$\bar t_1=3/2$, and $\bar t_i=3$ for $i=2,3,4$.
All four UAVs start idle at the depot with zero payload and full battery.
Their airspeeds are $(1,2,1,1)$; each has empty mass 1, payload capacity 1,
$\gamma_u=1/10$, battery capacity 10, and reserve 1.
All ground durations and energies are zero and bay capacity is unconstrained.

Each UAV has two active slots, with slot 2 committed to the terminal depot.
Commit the full transfer vectors: initial loading 1, first-slot delivery 1,
and all other transfers zero. The incumbent gives task $i$ to UAV $i$,
with a straight outbound path and its reverse. Assignments and first-slot
task identities are mutable. The positive first-slot delivery precludes
a depot or inactive first slot. Assignment consistency therefore makes the
four first-slot recipients a permutation of the four tasks.

Use two cells, east $\{(x,y)\in\mathcal W:x>0\}$ and west
$\{(x,y)\in\mathcal W:x\le0\}$, with one time bin $[0,4]$.
The previous wind is zero. The updated east wind is $(-1/2,0)$ and the
west wind remains zero. Thus the change affects a proper subset of
cell--time entries and its magnitude is strictly below every airspeed.
The incumbent is feasible under the previous wind, including all reserve
checks. Its eastbound times are initially 1 and $1/2$.

Under the updated wind the direct eastbound times of UAVs 1 and 2 are
2 and $2/3$, and the corresponding return times are $2/3$ and $2/5$.
UAV 1 misses recipient 1's deadline. Assigning recipient 1 to UAV 2 and
recipient 2 to UAV 1 satisfies both deadlines. The unaffected west UAVs
arrive at time 1 and return at time 2.
The exchanged plan has the following exact records.

\begin{table}[t]
\centering
\caption{Exact records of the constructed recovery witness.}
\label{tab:worked_recovery}
\begin{tabular}{ccccc}
\toprule
UAV & Task & Delivery time & Return time & Total energy\\
\midrule
1 & 2 & $2$ & $8/3$ & $7/15$\\
2 & 1 & $2/3$ & $16/15$ & $26/75$\\
3 & 3 & $1$ & $2$ & $3/10$\\
4 & 4 & $1$ & $2$ & $3/10$\\
\bottomrule
\end{tabular}
\end{table}

Outbound payload is one unit and return payload is zero. For example,
UAV 1 expends $(1/10)(1+1)(1)(2)=2/5$ outbound and
$(1/10)(1)(1)(2/3)=1/15$ on return. All terminal batteries exceed reserve,
and each post-service state uses that same direct return as its reserve
witness. The initial depot witness is the constant path. Total loading
is exactly four and all demand, payload, and horizon constraints hold.

\emph{Case A: structurally local reassignment.}
Commit both nominal path coordinates of every UAV to their incumbent
curves. Endpoint consistency then forces UAVs 1--2 to the east site and
UAVs 3--4 to the west site, while co-located recipient identities can swap.
Their potential participants are respectively $\{1,2\}$ and $\{3,4\}$.
Only first-slot sequence and assignment coordinates remain mutable.
The stock record is constant because all loads are committed; its mutable
support is empty. There is no bay constraint. For a per-task service
record, endpoints prove that only sequences in that task's group can
select it. All remaining argument records depend on at most one group,
so group-local supports satisfy the common-domain support condition.
Service constraints connect the two recipients within each group.

The paired diagnostic seeds recipients 1 and 2, so
$\mathcal V^R=\{1,2\}$. There are four structural permutations: the
identity or exchange within each group. The Level-2 space contains only
the incumbent, Level-3 also admits the east exchange, and Level-4 also
permits the west exchange. Thus
$\mathcal X_2\subsetneq\mathcal X_3\subsetneq\mathcal X_4$.
Only the east exchange restores recipient 1's deadline; hence
$k^F=k_{\mathfrak R}^F=3$. Every smaller level retains that recipient's
responsible slow UAV and its committed incoming path, so lower-level
infeasibility is exact, not a failed candidate search.

\emph{Case B: collapse under route participation.}
Make the nominal path coordinates mutable, retaining all other inputs
and commitments. Structurally every UAV can now serve every task;
$\mathcal T_u=\mathcal V^{\mathrm{rem}}$ for all $u$. A route-dependent
completion constraint therefore has full task incidence. The nonempty
east seed set gives $\mathcal V^R=M^A$ and
$\mathcal X_3=\mathcal X_4$. The same east exchange remains feasible.

In fact $k^F=3$ even with arbitrary admissible paths. Any route by UAV 1
from $x=0$ to $x=1$ must gain one unit of eastward displacement within
the east cell. There its eastward ground-velocity component is at most
$v_1-1/2=1/2$, so it needs at least two time units. Extra west-cell
excursions cannot reduce the required east-cell displacement. Since
$\bar t_1=3/2$, retaining UAV 1's responsibility cannot meet the positive
mandatory delivery at Levels 0--2. The reassignment witness proves the
matching upper bound.

For a finite reference library containing only the direct segment for
each endpoint pair, Case A enumerates four structural assignments and
Case B enumerates 24. Their feasible counts by level are respectively
$(0,0,0,1,2)$ and $(0,0,0,6,6)$; their pre-update feasible counts over
the full represented domain are four and 24. These are exact constructed
checks, not sampled mission results. The example demonstrates possible
strict locality under existing commitments and its loss when participation
widens. It does not establish that co-location or such commitments are
typical in disaster relief.
